\paragraph{Model scale and architecture.}
Our primary experiments span 7--14B parameters: DeepSeek-R1-Distill-Qwen-7B, Qwen2.5-7B-Instruct, and Qwen2.5-14B-Instruct (\S\ref{sec:scale-14b}). Whether the phase-modulated severity patterns persist, amplify, or attenuate at substantially larger scales (e.g., 70B) remains untested. Cross-architecture validation on DeepSeek-R1-Distill-Llama-8B (LlamaForCausalLM, \S\ref{sec:cross-arch}, Appendix~\ref{app:cross-arch}) supports that the phase gradient extends beyond the Qwen2 architecture ($n{=}250$, all 8 conditions significant; importance $d_z{=}0.81$--$1.15$), but strategy specificity is substantially weaker ($2.0$--$2.6{\times}$ vs.\ $15$--$17{\times}$ on DeepSeek) because random eviction also exhibits significant phase gradients on LLaMA ($d_z{=}0.41$--$0.46$). This architecture-dependent strategy specificity indicates that the degree of strategy${\times}$phase interaction, not just absolute magnitude, varies across attention implementations. The pattern has been tested on two architecture families (Qwen2ForCausalLM, LlamaForCausalLM) at two scale points (7B, 14B); comprehensive validation across more diverse architectures (e.g., Mistral, multi-latent attention) and larger scales is needed for full generalization claims.

\paragraph{Task coverage.}
We evaluate on three benchmarks spanning mathematical and scientific reasoning: MATH-500 (primary), GSM8K (\S\ref{sec:cross-benchmark}), and GPQA Diamond (\S\ref{sec:cross-benchmark-gpqa}). Cross-benchmark validation on GSM8K supports the phase gradient direction under importance-based eviction but reveals that random eviction is non-significant (CIs spanning zero), precluding meaningful ratio computation. On MATH-500, strategy specificity is quantifiable ($15$--$17{\times}$); on GSM8K, it manifests as a qualitative gap (significant vs.\ non-significant). This suggests that the phase-strategy interaction is robust in direction but task-dependent in form: complex reasoning (MATH-500) produces quantitative strategy contrast while simpler arithmetic (GSM8K) produces qualitative separation. Qwen GSM8K outputs were truncated at \texttt{max\_new\_tokens}${=}$2047, so Qwen GSM8K absolute values carry this limitation; DeepSeek GSM8K data is untruncated. Our task-complexity explanation for specificity attenuation (lower cognitive load $\to$ more uniform attention $\to$ reduced strategy sensitivity) is supported by within-benchmark correlational evidence but requires validation across a broader range of task complexities; whether different reasoning tasks (code generation, multi-hop QA) produce yet different attenuation profiles remains open. Random eviction on GSM8K exhibits numerical instability (extreme outlier II values), likely a systematic interaction between reduced-precision arithmetic and score-independent eviction.

\paragraph{Recency control architecture scope.}
The three-way dissociation (importance/random/recency) was demonstrated on the Qwen2 architecture. On LlamaForCausalLM with grouped-query attention, even random eviction produces statistically significant phase gradients ($d_z{=}0.41$--$0.46$), suggesting that GQA's reduced KV-head count may inherently amplify phase-correlated compound damage regardless of selection criterion. Recency-based validation on GQA architectures remains for future work.

\paragraph{Eviction strategy scope and H$_2$O confound.}
Our importance-based eviction uses cumulative attention scores (H$_2$O-style scoring~\cite{h2o-neurips2023}, Eq.~\ref{eq:kv-importance}), which creates an inherent confound: head pruning alters the attention distribution from which importance scores are computed, and eviction then operates on these degraded scores. The compound interaction patterns we report, including importance discrimination collapse and the phase-strategy crossover, may be partly attributable to this circular dependency rather than to compound sparsity \emph{per se}. Other attention-mass-based methods (SnapKV~\cite{snapkv-neurips2024}, PyramidKV~\cite{pyramidkv-colm2025}) share the same scoring basis and would likely exhibit a similar confound. The recency control (\S\ref{sec:scoring-control}) partially addresses this by confirming that the phase gradient requires the attention-score pathway ($d_z{=}0.083$, NS under recency eviction vs.\ $d_z{=}0.574$ under importance at matched $s{=}0.15$). However, recency eviction was tested at a single sparsity level ($n{=}100$ pilot, vs.\ $n{=}500$ for importance); comprehensive evaluation across sparsity levels, and testing of other score-independent strategies (e.g., learned eviction policies), remains an important next step. The qualitative conclusion (gradient absence, CI spanning zero) is robust to sample size: post-hoc power analysis indicates $n{=}100$ provides power $> 0.85$ to detect effects of $d_z \geq 0.30$ at $\alpha{=}0.05$, well below the observed importance effect ($d_z{=}0.57$).

\paragraph{Symmetric sparsity assumption.}
Our analysis applies symmetric sparsity levels: the same fraction $s$ is used for both head pruning and KV cache eviction. A preliminary spot check under asymmetric allocation (Appendix~\ref{app:asymmetric-robustness}) supports that the phase gradient is robust to head/KV budget ratios, but systematic characterization of the full asymmetric interaction surface remains future work.

\paragraph{Phase segmentation validation.}
The \texttt{<think>} token validation of sparsity-based phase segmentation (\S\ref{sec:background}, Appendix~\ref{app:think-validation}) is limited to a single model (DeepSeek-R1-Distill-Qwen-7B) on a single benchmark (GSM8K, $n{=}48$). Qwen2.5-7B-Instruct does not generate explicit reasoning markers, precluding the same validation. Whether the high agreement ($\kappa{=}0.88$) generalizes to other models with explicit CoT delimiters, or to non-mathematical reasoning tasks, remains open.

\paragraph{Mechanism.}
Importance discrimination collapse is a correlational diagnostic observation, not a causal explanation: both strategies exhibit near-identical Gini drops yet severity diverges by $15$--$17{\times}$, indicating that Gini collapse alone is insufficient. Causal validation through controlled entropy manipulation is needed. The entropy redistribution component replicates on only 2/4 Qwen conditions, suggesting model-specific factors.

\paragraph{Perplexity vs.\ accuracy.}
Our primary severity metric is perplexity-based (II), chosen because its per-token granularity enables the phase-level decomposition central to our analysis; accuracy is binary per-problem and cannot be meaningfully decomposed by reasoning phase. Accuracy validation on MATH-500 (\S\ref{sec:accuracy-validation}, Appendix~\ref{app:accuracy-details}) shows gradual decline through $s{=}0.25$ ($\delta{=}-4$ to $-13$\,pp) followed by a sharp cliff at $s{=}0.30$ ($\delta{=}-27$\,pp) for Qwen ($n{=}100$, truncation $\leq 4\%$). DeepSeek accuracy deltas are confounded by rising truncation ($16\%$ to $76\%$); non-truncated analysis reveals a dual-regime structure but is limited by small surviving samples at high sparsity. Comprehensive accuracy evaluation across benchmarks remains an important next step.

\paragraph{Sparsity range.}
Our experiments span $s \in \{0.15, 0.20, 0.25, 0.30\}$, sufficient to establish monotonicity at the tested levels (JT test $p \ll 10^{-10}$) but insufficient to characterize the functional form of the auxiliary ratio $G(s)$ or identify potential saturation thresholds. More extreme sparsity ($s > 0.50$) may exhibit qualitatively different interaction dynamics. We do not report inference throughput or latency measurements, which are necessary for a complete deployment analysis but orthogonal to our severity characterization. Extended limitations in Appendix~\ref{app:limitations}.
